
\documentclass{article}
\usepackage{colortbl}
\usepackage{makecell}
\usepackage{multirow}
\usepackage{supertabular}

\begin{document}

\newcounter{utterance}

\centering \large Interaction Transcript for game `ifeval', experiment `ifeval', episode 303 with Qwen3.5{-}4B.
\vspace{24pt}

{ \footnotesize  \setcounter{utterance}{1}
\setlength{\tabcolsep}{0pt}
\begin{supertabular}{c@{$\;$}|p{.15\linewidth}@{}p{.15\linewidth}p{.15\linewidth}p{.15\linewidth}p{.15\linewidth}p{.15\linewidth}}
    \# & \multicolumn{2}{c}{Player} && \multicolumn{2}{c}{Game Master} \\
    \hline

    \theutterance \stepcounter{utterance}  
    & & & \multicolumn{4}{p{0.6\linewidth}}{
        \cellcolor[rgb]{0.9,0.9,0.9}{
            \makecell[{{p{\linewidth}}}]{
                \texttt{\tiny{[P1$\langle$GM]}}
                \texttt{Write a paragraph that lists the average length of various animal specimens from smallest to largest. Your response should contain less than 17 sentences.} \\
            }
        }
    }
     \\ \\

    \theutterance \stepcounter{utterance}  
    & & \multicolumn{4}{p{0.6\linewidth}}{
        \cellcolor[rgb]{0.95,0.95,0.95}{
            \makecell[{{p{\linewidth}}}]{
                \texttt{\tiny{[P1$\leftrightarrow$P]}}
                \texttt{{[}DRAFT{]}} \\
\texttt{The average length of animal specimens varies wildly across the spectrum of life on Earth. At the very bottom of this scale are microscopic organisms like the *Paramecium caudatum*, which measures only about 0.05 millimeters. Moving slightly larger, we find the *Artemia salina*, a brine shrimp that typically grows to roughly 2 millimeters. Among insects, the *Gonolabis* mosquito is notably small, averaging just 1.5 millimeters in length. The *Balanoglossus* acorn worm represents a larger invertebrate, reaching approximately 10 centimeters. In the vertebrate world, the *Microtus* vole is quite tiny, with an average body length of about 12 centimeters. The *Peromyscus* deer mouse is slightly bigger, standing around 15 centimeters long. Larger mammals include the *Mustela* weasel, which averages 40 centimeters, and the *Vulpes* fox, which reaches about 60 centimeters. The *Canis* wolf is significantly larger, with an average length of 1.5 meters. At the top of this specific list of relatively small animals is the *Oryctolagus* rabbit, which can grow to approximately 1 meter. While there are many species larger than the rabbit, such as elephants and blue whales, this selection illustrates the progression from microscopic to macroscopic forms within a manageable range. Each of these specimens represents a unique evolutionary adaptation to its specific environment.} \\
            }
        }
    }
    & \\ \\

    \theutterance \stepcounter{utterance}  
    & & \multicolumn{4}{p{0.6\linewidth}}{
        \cellcolor[rgb]{0.95,0.95,0.95}{
            \makecell[{{p{\linewidth}}}]{
                \texttt{\tiny{[P1$\leftrightarrow$P]}}
                \texttt{{[}SELF{-}FEEDBACK{]}} \\
\texttt{FORMAT: ok} \\
\texttt{RULES: problem} \\
\texttt{STRATEGY: problem} \\
\texttt{REASON: The answer fails to follow the specific instruction to list specimens from smallest to largest because it includes the rabbit (1 meter) after the wolf (1.5 meters), violating the ordering constraint.} \\
            }
        }
    }
    & \\ \\

    \theutterance \stepcounter{utterance}  
    & \multicolumn{4}{p{0.6\linewidth}}{
        \cellcolor[rgb]{0.561,0.737,0.561}{
            \makecell[{{p{\linewidth}}}]{
                \texttt{\tiny{[P1$\rangle$GM]}}
                \texttt{FORMAT: ok} \\
\texttt{RULES: The answer fails to follow the specific instruction to list specimens from smallest to largest because it includes the rabbit (1 meter) after the wolf (1.5 meters), violating the ordering constraint.} \\
\texttt{STRATEGY: Reorder the listed animal specimens strictly by their average length, ensuring the sequence progresses continuously from the smallest organism to the largest without any backtracking or out{-}of{-}order entries.} \\
\texttt{REASON: The original text placed the rabbit after the wolf, which breaks the required ascending order of size. The refined version must arrange all mentioned animals so that each subsequent specimen is larger than the previous one, maintaining a logical and accurate progression from the microscopic Paramecium to the largest mammal included in the list.} \\
            }
        }
    }
    & & \\ \\

    \theutterance \stepcounter{utterance}  
    & & & \multicolumn{2}{p{0.3\linewidth}}{
        \cellcolor[rgb]{0.9,0.9,0.9}{
            \makecell[{{p{\linewidth}}}]{
                \texttt{\tiny{[GM$|$GM]}}
                {'length_constraints:number_sentences': {'relation': 'less than', 'num_sentences': 17}}
            }
        }
    }
    & & \\ \\

    \theutterance \stepcounter{utterance}  
    & & & \multicolumn{2}{p{0.3\linewidth}}{
        \cellcolor[rgb]{0.9,0.9,0.9}{
            \makecell[{{p{\linewidth}}}]{
                \texttt{\tiny{[GM$|$GM]}}
                \texttt{game\_result = WIN} \\
            }
        }
    }
    & & \\ \\

\end{supertabular}
}

\end{document}
